% GENERATED FILE. Edit canonical lessons, then run npm run generate:books.
% canonical-source-hash: fnv1a64:d1d7a5a437ea602d
% canonical-lessons: MR-C388-pradhanya, MR-C388-uddista, MR-C388-pragti, MR-C388-vikas, MR-C388-amlat-anne, MR-R388-review-382-388

\chapter{Goals and Progress}
\label{ch:mr-goals-and-progress}
\hlchaptermodality{388}

\begin{chapteropening}
\textbf{By the end of this chapter:} \emph{I can talk about priorities, targets and progress.}

Complete the last lesson of chapter 388: I can talk about priorities, targets and progress.
\end{chapteropening}

\section[\texorpdfstring{\mr{प्राधान्य} (prādhānya)}{प्राधान्य (prādhānya)}]{\mr{प्राधान्य} (prādhānya) — a priority}

\label{lesson:MR-C388-pradhanya}



\begin{quote}
Before the new one: say the Marathi for to turn out, then the Marathi for to be convinced.
\end{quote}

\subsection*{You'll want to know: \mr{प्राधान्य}}
\textbf{\mr{प्राधान्य}} — \emph{prādhānya} — \textquotedblleft{}a priority\textquotedblright{}.

\begin{quote}
\mr{माझ्यासाठी} \mr{आरोग्याला} \mr{प्राधान्य} \mr{आहे}.
\end{quote}

\begin{quote}
\emph{mājhyāsāṭhī ārogyālā prādhānya āhe.}
\end{quote}

\begin{quote}
For me, health comes first.
\end{quote}

Neuter.

X \mr{ला} \mr{प्राधान्य} \mr{देणे}, to give priority to X.

\begin{cousinweb}
\textbf{\mr{प्राधान्य}} (\emph{prādhānya}) is a tatsama, taken from Sanskrit \emph{prādhānya}, \textquotedblleft{}pre-eminence, predominance\textquotedblright{} (Monier-Williams).
\end{cousinweb}

\begin{grammarlens}[title={What you've built}]
One more word for saying what happens if.
\end{grammarlens}

\subsection*{Your turn}
\begin{itemize}
\raggedright
  \item \emph{Say it:} \emph{prādhānya}
  \item \emph{Say it:} \emph{prādhānya}, once more
  \item \emph{Say it:} \emph{mājhyāsāṭhī ārogyālā prādhānya āhe.} — For me, health comes first.
  \item \emph{Recall:} say the Marathi for to turn out, then the Marathi for to be convinced, then say \emph{prādhānya} again
\end{itemize}

\begin{tcolorbox}[breakable,colback=teal!4,colframe=teal!35!black,title={Before you move on}]
What does \mr{प्राधान्य} mean? (\textquotedblleft{}A priority\textquotedblright{}.) Say it once more.
\end{tcolorbox}
\section[\texorpdfstring{\mr{उद्दिष्ट} (uddiṣṭa)}{उद्दिष्ट (uddiṣṭa)}]{\mr{उद्दिष्ट} (uddiṣṭa) — an objective; a target}

\label{lesson:MR-C388-uddista}



\begin{quote}
Before the new one: say the Marathi for to be convinced, then the Marathi for a priority.
\end{quote}

\subsection*{You'll want to know: \mr{उद्दिष्ट}}
\textbf{\mr{उद्दिष्ट}} — \emph{uddiṣṭa} — \textquotedblleft{}an objective; a target\textquotedblright{}.

\begin{quote}
\mr{या} \mr{वर्षी} \mr{शंभर} \mr{झाडं} \mr{लावण्याचं} \mr{उद्दिष्ट} \mr{आहे}.
\end{quote}

\begin{quote}
\emph{yā varṣī śambhar jhāḍaṁ lāvaṇyācaṁ uddiṣṭa āhe.}
\end{quote}

\begin{quote}
The objective this year is to plant a hundred trees.
\end{quote}

Neuter.

Neuter: \mr{उद्दिष्ट} \mr{गाठणे}, to reach a target.

\begin{cousinweb}
\textbf{\mr{उद्दिष्ट}} (\emph{uddiṣṭa}) is a tatsama, taken from Sanskrit \emph{uddiṣṭa}, \textquotedblleft{}mentioned, specified\textquotedblright{} (Monier-Williams): what has been named as the aim.
\end{cousinweb}

\begin{grammarlens}[title={What you've built}]
One more word for saying what happens if.
\end{grammarlens}

\subsection*{Your turn}
\begin{itemize}
\raggedright
  \item \emph{Say it:} \emph{uddiṣṭa}
  \item \emph{Say it:} \emph{uddiṣṭa}, once more
  \item \emph{Say it:} \emph{yā varṣī śambhar jhāḍaṁ lāvaṇyācaṁ uddiṣṭa āhe.} — The objective this year is to plant a hundred trees.
  \item \emph{Recall:} say the Marathi for to be convinced, then the Marathi for a priority, then say \emph{uddiṣṭa} again
\end{itemize}

\begin{tcolorbox}[breakable,colback=teal!4,colframe=teal!35!black,title={Before you move on}]
What does \mr{उद्दिष्ट} mean? (\textquotedblleft{}An objective; a target\textquotedblright{}.) Say it once more.
\end{tcolorbox}
\section[\texorpdfstring{\mr{प्रगती} (pragtī)}{प्रगती (pragtī)}]{\mr{प्रगती} (pragtī) — progress}

\label{lesson:MR-C388-pragti}



\begin{quote}
Before the new one: say the Marathi for a priority, then the Marathi for an objective.
\end{quote}

\subsection*{You'll want to know: \mr{प्रगती}}
\textbf{\mr{प्रगती}} — \emph{pragtī} — \textquotedblleft{}progress\textquotedblright{}.

\begin{quote}
\mr{या} \mr{वर्षी} \mr{अरुणने} \mr{अभ्यासात} \mr{चांगली} \mr{प्रगती} \mr{केली}.
\end{quote}

\begin{quote}
\emph{yā varṣī aruṇne abhyāsāt cāṅglī pragtī kelī.}
\end{quote}

\begin{quote}
Arun made good progress in his studies this year.
\end{quote}

Feminine.

\mr{प्रगती} \mr{करणे}, to make progress.

\begin{cousinweb}
\textbf{\mr{प्रगती}} (\emph{pragtī}) is a learned word built from Sanskrit \emph{pra-}, \textquotedblleft{}forward\textquotedblright{}, and \emph{gati}, \textquotedblleft{}going\textquotedblright{} (Monier-Williams).
\end{cousinweb}

\begin{grammarlens}[title={What you've built}]
One more word for saying what happens if.
\end{grammarlens}

\subsection*{Your turn}
\begin{itemize}
\raggedright
  \item \emph{Say it:} \emph{pragtī}
  \item \emph{Say it:} \emph{pragtī}, once more
  \item \emph{Say it:} \emph{yā varṣī aruṇne abhyāsāt cāṅglī pragtī kelī.} — Arun made good progress in his studies this year.
  \item \emph{Recall:} say the Marathi for a priority, then the Marathi for an objective, then say \emph{pragtī} again
\end{itemize}

\begin{tcolorbox}[breakable,colback=teal!4,colframe=teal!35!black,title={Before you move on}]
What does \mr{प्रगती} mean? (\textquotedblleft{}Progress\textquotedblright{}.) Say it once more.
\end{tcolorbox}
\section[\texorpdfstring{\mr{विकास} (vikās)}{विकास (vikās)}]{\mr{विकास} (vikās) — development}

\label{lesson:MR-C388-vikas}



\begin{quote}
Before the new one: say the Marathi for an objective, then the Marathi for progress.
\end{quote}

\subsection*{You'll want to know: \mr{विकास}}
\textbf{\mr{विकास}} — \emph{vikās} — \textquotedblleft{}development\textquotedblright{}.

\begin{quote}
\mr{गावाच्या} \mr{विकासासाठी} \mr{नवा} \mr{रस्ता} \mr{बांधला}.
\end{quote}

\begin{quote}
\emph{gāvācyā vikāsāsāṭhī navā rastā bāndhlā.}
\end{quote}

\begin{quote}
A new road was built for the development of the village.
\end{quote}

Masculine.

X \mr{चा} \mr{विकास}, the development of X.

\begin{cousinweb}
\textbf{\mr{विकास}} (\emph{vikās}) is a tatsama, taken from Sanskrit \emph{vikāsa}, \textquotedblleft{}expanding, flowering, growth\textquotedblright{} (CDIAL 11630).
\end{cousinweb}

\begin{grammarlens}[title={What you've built}]
One more word for saying what happens if.
\end{grammarlens}

\subsection*{Your turn}
\begin{itemize}
\raggedright
  \item \emph{Say it:} \emph{vikās}
  \item \emph{Say it:} \emph{vikās}, once more
  \item \emph{Say it:} \emph{gāvācyā vikāsāsāṭhī navā rastā bāndhlā.} — A new road was built for the development of the village.
  \item \emph{Recall:} say the Marathi for an objective, then the Marathi for progress, then say \emph{vikās} again
\end{itemize}

\begin{tcolorbox}[breakable,colback=teal!4,colframe=teal!35!black,title={Before you move on}]
What does \mr{विकास} mean? (\textquotedblleft{}Development\textquotedblright{}.) Say it once more.
\end{tcolorbox}
\section[\texorpdfstring{\mr{अमलात} \mr{आणणे} (amlāt āṇṇe)}{अमलात आणणे (amlāt āṇṇe)}]{\mr{अमलात} \mr{आणणे} (amlāt āṇṇe) — to implement; to put into effect}

\label{lesson:MR-C388-amlat-anne}



\begin{quote}
Before the new one: say the Marathi for progress, then the Marathi for development.
\end{quote}

\subsection*{You'll want to know: \mr{अमलात} \mr{आणणे}}
\textbf{\mr{अमलात} \mr{आणणे}} — \emph{amlāt āṇṇe} — \textquotedblleft{}to implement; to put into effect\textquotedblright{}.

\begin{quote}
\mr{समितीने} \mr{ठरलेली} \mr{योजना} \mr{लगेच} \mr{अमलात} \mr{आणली}.
\end{quote}

\begin{quote}
\emph{samitīne ṭharlelī yojnā lagec amlāt āṇlī.}
\end{quote}

\begin{quote}
The committee put the agreed plan into effect at once.
\end{quote}

Past, with -\mr{ने} (he …): \mr{त्याने} \mr{अमलात} \mr{आणलं}, \emph{tyāne amlāt āṇlaṁ}.

\mr{अमलात} \mr{येणे}, to come into effect.

\begin{cousinweb}
\textbf{\mr{अमलात} \mr{आणणे}} (\emph{amlāt āṇṇe}) joins \mr{अमलात}, \textquotedblleft{}into effect\textquotedblright{}, from \mr{अमल}, a loan from Persian usually traced to Arabic \emph{ʿamal}, \textquotedblleft{}action\textquotedblright{}, to \mr{आणणे}, \textquotedblleft{}to bring\textquotedblright{}.
\end{cousinweb}

\begin{grammarlens}[title={What you've built}]
One more word for saying what happens if.
\end{grammarlens}

\subsection*{Your turn}
\begin{itemize}
\raggedright
  \item \emph{Say it:} \emph{amlāt āṇṇe}
  \item \emph{Say it:} \emph{amlāt āṇṇe}, once more
  \item \emph{Say it:} \emph{samitīne ṭharlelī yojnā lagec amlāt āṇlī.} — The committee put the agreed plan into effect at once.
  \item \emph{Recall:} say the Marathi for progress, then the Marathi for development, then say \emph{amlāt āṇṇe} again
\end{itemize}

\begin{tcolorbox}[breakable,colback=teal!4,colframe=teal!35!black,title={Before you move on}]
What does \mr{अमलात} \mr{आणणे} mean? (\textquotedblleft{}To implement; to put into effect\textquotedblright{}.) Say it once more.
\end{tcolorbox}
\section[(dialogue)]{Looking back: from Likely or Certain to Goals and Progress}

\label{lesson:MR-R388-review-382-388}



\begin{quote}
Say the Marathi for likely and probably.
\end{quote}

\subsection*{Your turn}
Say these aloud:

\begin{itemize}
\raggedright
  \item likely, probably, naturally, certain, to seem
  \item a prediction, to assume, an expectation, a trend, to size up
  \item circumstances, a state, a plight, a tight spot, a test
  \item to worsen, to surge up, to calm down, to be averted, to confront
  \item to prevent, to restrain, to block, to take notice, a loophole
  \item to leave out, to include, to rely on, to turn out, to be convinced
  \item a priority, an objective, progress, development, to implement
\end{itemize}

\begin{tcolorbox}[breakable,colback=teal!4,colframe=teal!35!black,title={Before you move on}]
Likely? (\textbf{\mr{संभाव्य}}, \emph{sambhāvya}.) Probably? (\textbf{\mr{बहुधा}}, \emph{bahudhā}.) Naturally? (\textbf{\mr{साहजिकच}}, \emph{sāhjikac}.) Certain? (\textbf{\mr{निश्चित}}, \emph{niścit}.) To seem? (\textbf{\mr{भासणे}}, \emph{bhāsṇe}.) A prediction? (\textbf{\mr{भाकीत}}, \emph{bhākīt}.) To assume? (\textbf{\mr{गृहीत} \mr{धरणे}}, \emph{gṛhīt dharṇe}.) An expectation? (\textbf{\mr{अपेक्षा}}, \emph{apekṣā}.) A trend? (\textbf{\mr{कल}}, \emph{kal}.) To size up? (\textbf{\mr{ताडणे}}, \emph{tāḍṇe}.) Circumstances? (\textbf{\mr{परिस्थिती}}, \emph{paristhitī}.) A state? (\textbf{\mr{स्थिती}}, \emph{sthitī}.) A plight? (\textbf{\mr{अवस्था}}, \emph{avasthā}.) A tight spot? (\textbf{\mr{पेच}}, \emph{pec}.) A test? (\textbf{\mr{कसोटी}}, \emph{kasoṭī}.) To worsen? (\textbf{\mr{चिघळणे}}, \emph{cighaḷṇe}.) To surge up? (\textbf{\mr{उफाळणे}}, \emph{uphāḷṇe}.) To calm down? (\textbf{\mr{निवळणे}}, \emph{nivaḷṇe}.) To be averted? (\textbf{\mr{टळणे}}, \emph{ṭaḷṇe}.) To confront? (\textbf{\mr{सामोरे} \mr{जाणे}}, \emph{sāmore jāṇe}.) To prevent? (\textbf{\mr{प्रतिबंध} \mr{करणे}}, \emph{pratibandh karṇe}.) To restrain? (\textbf{\mr{रोखणे}}, \emph{rokhṇe}.) To block? (\textbf{\mr{अडवणे}}, \emph{aḍavṇe}.) To take notice? (\textbf{\mr{दखल} \mr{घेणे}}, \emph{dakhal gheṇe}.) A loophole? (\textbf{\mr{पळवाट}}, \emph{paḷvāṭ}.) To leave out? (\textbf{\mr{वगळणे}}, \emph{vagaḷṇe}.) To include? (\textbf{\mr{समाविष्ट} \mr{करणे}}, \emph{samāviṣṭa karṇe}.) To rely on? (\textbf{\mr{विसंबणे}}, \emph{visambṇe}.) To turn out? (\textbf{\mr{ठरणे}}, \emph{ṭharṇe}.) To be convinced? (\textbf{\mr{पटणे}}, \emph{paṭṇe}.) A priority? (\textbf{\mr{प्राधान्य}}, \emph{prādhānya}.) An objective? (\textbf{\mr{उद्दिष्ट}}, \emph{uddiṣṭa}.) Progress? (\textbf{\mr{प्रगती}}, \emph{pragtī}.) Development? (\textbf{\mr{विकास}}, \emph{vikās}.) To implement? (\textbf{\mr{अमलात} \mr{आणणे}}, \emph{amlāt āṇṇe}.)
\end{tcolorbox}
